\cotefile{64}
\scanpage{91}
les \deleted{\unsure{objets|objet}} épinglages de Legendre-Jacobi ``directs'', i.e. tels que la section de $\mu_4^\times$ qu'\deleted{\unsure{elles}}il définit \unsure{soit|sont} justement $i$. \deleted{\unsure{leurs|Leurs|base}} \deleted{Les}Le \inserted{champ des} courbes elliptiques relatives \inserted{(sur des $S/\widetilde S_0$)} munies d'un épinglage de Legendre-Jacobi direct, est représentable par $U_{03\,\widetilde S_0}$, sur lequel on fait opérer le groupe modulaire $\mathbb D'_3$, qui est ainsi \inserted{justement} le sous-groupe de $\mathbb D_3\times\mu_{4\,\widetilde S_0}$ qui fait passer d'un épinglage direct à un épinglage direct.

Revenant au cas de la base absolue $\unsure{S_0|L_0}$, on doit encore expliciter les opérations de $\mu_4\times\unsure{\mathfrak S_3|\mathcal G_3}$ sur $\underline O_{U_{0,3}\times\mu_4^\times}\deleted{\unsure{(-1)|(1)}}(1)$, via l'action tautologique résultant de l'action de $\mathfrak S_3$ sur $V_{\{0,1,\infty\}}(\underline O_{S_0})$ d'une part, de \inserted{$\mathfrak S_3$ et} $\mu_4$ sur le facteur $\mu_4^\times$ d'autre part. On désigne l'action tautologique d'un $(g,\lambda)$ sur $\zeta\in\Gamma\underline O_{\unsure{U_{0,3}|\mathbb P_{0,3}|U_{03}}}(1)(\unsure{1|t|A})$
% archive sheet number: 90
\scanpage{92}
\mspage{10}
\inserted{(où $t:S\longrightarrow U_{03}$)} par la notation \unsure{anodine|ordinaire}
\[
(g,\lambda).\zeta\;,
\]
et on trouve que l'on a, pour l'action modulaire $T_{g,\lambda}$ sur $\zeta$ :
\[
T_{g,\lambda}\zeta=\lambda\deleted{\operatorname{sg}(g)}(g,\lambda).\zeta
\qquad \unsure{\zeta}\in\Gamma\underline O_{\unsure{U_{0,3}|\mathbb P_{0,3}|U_{03}}}(1)(t) \tag*{(51)}
\]
donc
\[
T_{g,\lambda}\deleted{\unsure{\zeta^2|\zeta}}y=\lambda^2(g,\lambda).\deleted{\unsure{\zeta^2}}y
\qquad y\in\Gamma\underline O_{U_{0,3}}(2)(t) \tag*{(52)}
\]
donc
\[
\begin{aligned}
T_{g,\lambda}(\tau)&=\lambda^2\operatorname{sg}(g)\,(g,\lambda).\tau
\qquad \tau\in\Gamma\deleted{\unsure{\underline O_{U_{03}}(t)}}\\
&\hphantom{{}=\lambda^2\operatorname{sg}(g)\,(g,\lambda).\tau\qquad \tau\in\Gamma}
\underline O_{U_{0,3}}(2)\unsure{(\tfrac12)_0|(\tfrac12)}(t)
\end{aligned} \tag*{(53)}
\]

Quand on se restreint au sous-groupe $\mathbb D'_{3\,S_0}$, pour lequel on a $\lambda^2=\operatorname{sg}(g)$, on trouve donc pour (51) (52) (53)
\[
T_{g\deleted{\unsure{,\lambda}}}\deleted{\lambda}\deleted{\unsure{y|\zeta}}\zeta=
\deleted{\operatorname{sg}(g)}\chi_4(g)\,(g.\zeta) \tag*{(51 bis)}
\]
\[
T_g\,y=\operatorname{sg}(g^0)\,(g.y) \tag*{(52 bis)}
\]
\[
T_g(\tau)=\deleted{\unsure{\lambda|1}}g.\tau \tag*{(53 bis)}
\]
où $g\in\Gamma(\mathbb D'_{3\,S})$ pour $S$ sur $S_0$
\[
\chi_4:\mathbb D'_{3\,S_0}\longrightarrow\mu_4
\quad\text{le caractère canonique} \tag*{(54)}
\]
(de sorte que $\mathbb D'_{3\,\mathrm{ab}}\simeq\mu_4$)
% archive sheet number: 91
\scanpage{93}
Ces dernières formules sont surtout intéressantes dans le contexte où on travaille \deleted{\unsure{la|ici}} sur \deleted{les}la base absolue $\widetilde S_0$ et avec les structures de Legendre\inserted{-Jacobi} directes, et où le schéma modulaire $U_{03}\times_{S_0}\mu_4^\times$ sur $S_0$ est remplacé par $U_{03\,\widetilde S_0}$ $(=U_{03}\times_{S_0}\mu_4^\times\,!)$ sur $\widetilde S_0$, et le groupe modulaire $\mathbb D_{3\,S_0}\times_{S_0}\mu_{4\,S_0}$ est remplacé par le \inserted{ss-}groupe $\mathbb D'_{3\,\widetilde S_0}$ de \deleted{\unsure{\mathbb D_3|\mathbb D_8}} son image inverse sur $\widetilde S_0$.

\begin{center}--\quad --\quad --\end{center}

Il faut maintenant définir les épinglages de Jacobi d'échelon $4$ relatifs sur l'échelon $2$, \deleted{i.e. les sections} \inserted{qui sont des sections \inserted{(\unsure{variables|véritables})}} de $ {}_4E\longrightarrow {}_2E$ au dessus de $\unsure{\underline I_E|\underline J_E}=({}_2E)^*$, \unsure{satisfont|satisfant|satisfaisant} la condition exprimée \deleted{\unsure{que|par}} sur chaque fibre géométrique par $\unsure{\zeta_i+\zeta_j+\zeta_k|\zeta_i+\zeta_j+\zeta_h|\xi_i+\xi_j+\xi_k}=0$, et montrer comment un tel épinglage définit canoniquement un épinglage de Legendre.
% archive sheet number: 92
\scanpage{94}
\mspage{11}
Mais avant, un commentaire. Quand on donne $\underline I$, rev. cubique étale de $S$, et une section $t$ de $\Sigma_{\underline I}^*$, alors la donnée d'une courbe elliptique relative sur $S$, correspondante à ces données $\unsure{\underline I|\underline\Pi|\underline J},t$, équivaut à celle d'un module inv. $T$ sur $S$, et d'un isom $(5)$
\[
T^{\otimes2}\simeq T_{\Sigma,t} \tag*{(*)}
\]
Quand on a dispose déjà d'un tel $T$, soit $T_0$, alors la catégorie des $T$ \inserted{sur $S$} munis de $(*)$ équivaut à celle des $\mu_{2\,S}$-torseurs $\underline\varepsilon$ sur $S$, moyennant l'équivalence
\[
\deleted{\unsure{\underline T}}\underline\varepsilon\longmapsto T=T_0\mathbin{\wedge}_{\mu_{2\,S}}\underline\varepsilon \tag*{(55)}
\]
Ici, on a envie de prendre
\[
T_0=\underline O_\Sigma(1)(t)^{\unsure{\sharp|\xi|\natural}}=\underline O_\Sigma(1)(t)\mathbin{\wedge}_{\mu_4}Q_{0S} \tag*{(56)}
\]
mais \deleted{\unsure{p}}on a un iso. can.
\[
T_0^{\otimes2}\simeq\underline O_\Sigma(2)(t)\mathbin{\wedge}_{\mu_2}\mu_4^*
\simeq T_{\Sigma,t}\mathbin{\wedge}_{\mu_2}(\mu_4^\times\mathbin{\wedge}\underline\omega_{\underline I})
\]
et pour avoir un iso. ``raisonnable'' $T_0^{\otimes2}\simeq T_{\Sigma,t}$, il faut se donner une \deleted{section} iso
\[
\alpha\unsure{\in|:}\deleted{\unsure{\sharp}}\mu_4^\times\simeq\underline\omega_{\underline I}
\quad\text{i.e. une section de }\mu_4^\times\mathbin{\wedge}\underline\omega_{\underline I}
\]
(ce qui est une des données indiquées par un épinglage de Legendre d'une courbe elliptique $/S$ correspondant à $\underline I$ \ldots). Une fois $\alpha$ fixé, on
% archive sheet number: 93
\scanpage{95}
trouve l'isom
\[
\unsure{\kappa_\alpha|K_\alpha|k_\alpha}:T_0^{\otimes2}\simeq T_{\Sigma,t}
\]
\unsure{d'où|où} une courbe elliptique $E_\alpha$, \inserted{corr. à $\underline I,t$} munie d'un \deleted{\unsure{rigidification}} \inserted{épinglage} de Legendre subordonné à $\alpha$ (puisque $T_{E_\alpha,0}\simeq T_0$ \ldots).
% slanted
\margin{c'est la courbe elliptique \uline{universelle}, correspondant à $\underline I,t$ et munie d'un épinglage de Legendre subordonné à $\alpha$ \ldots}
Ceci posé, pour une courbe elliptique quelconque $E$ sur $S$, corr. à $\underline I,t$, les épinglages de Legendre sur $E$ subordonnés à $\alpha$ forment (on l'a vu) les sections d'un $\mu_2$-torseur $\mathscr L_\alpha(E)$ sur $S$, et on voit de suite qu'on a un isom de $\mu_2$-torseurs
\[
\deleted{\mathscr L_\alpha(E)\longleftarrow}
\]
\[
\begin{aligned}
\underbrace{\underline{\operatorname{Isom}}_{S,\underline I}(E_\alpha,E)}_{
\substack{\text{l'indice }\underline I\text{ signifie}\\
\text{qu'on prend les iso.}\\
\text{compatibles avec}\\
\underline I\simeq {}_2E\simeq {}_2E_\alpha}}
&\xrightarrow{\sim}\mathscr L_\alpha(E)\\
u&\longmapsto u(\kappa_\alpha)
\end{aligned}\tag*{(57)}
\]
Si, au lieu d'une section \inserted{$\alpha$} de $\mu_4^\times\mathbin{\wedge}\underline\omega_{\underline I}$, on disposait d'une section $\omega$ de $\underline\omega_{\underline I}$ lui-même, on pourrait encore définir une courbe elliptique $E_\omega$ canonique corr. à $\underline I,t$, et munie d'un ``ép. de Legendre ``naïf'''' \inserted{(\unsure{i.e.|lire} d'un iso. $T_{E,0}\simeq\underline O_\Sigma(1)(t)$ tel que \ldots)}, et
% archive sheet number: 94
\scanpage{96}
jouant un rôle universel pour \deleted{\unsure{ces}} \inserted{de telles} données ($\underline I$ torseur sous $(\mathbb Z/3\mathbb Z)_S$, $t\in\Gamma(\Sigma^*_{\underline I}/S)$, $\omega\in\Gamma(S,\underline\omega_{\underline I})$ correspondant à la structure de torseur sur $\underline I$, et donnant lieu à une trivialisation de $\underline\omega_{\underline I}$ sur $S$ \ldots). Les deux pts de vue sont équivalents, quand on se fixe une section $i$ de $\mu_4^\times$, i.e. quand on travaille au dessus de la base $\widetilde S_0=\mu_{4\,S_0}^*$ ($S_0=\operatorname{Spec}\mathbb Z[1/2]$).

\begin{center}---\end{center}

Explicitons entièrement la notion d'épinglage de Legendre, du point de vue calculatoire, sur un corps de base alg. clos $k$ pour simplifier, en décrivant la courbe elliptique $E$ sur $k$ en termes de
\[
\boxed{I\bigl(={}_2E(k)^*\bigr)}\quad\text{(ens à trois éléments), du} \tag*{(58)}
\]
$k$-vectoriel
\[
\boxed{V_I=\deleted{\unsure{k^I|k}}\operatorname{Ker}(k^I\xrightarrow{\text{somme}}k)} \tag*{(59)}
\]
(\unsure{de là|et là|d'où} $\Sigma_I=\mathbb P^1\deleted{\unsure{roj}}(V_I)$, avec \deleted{\unsure{q|\delta}} des $\dot\xi_i\in\Sigma_I(k)$ $(i\in I)$ d'où $\Sigma_I^*=\Sigma_I-\{\dot\xi_i\}_{i\in I}$), et un $t\in\Sigma_I^\times$ i.e. une droite
% slanted
\margin{$\dot\xi_i=k.(e_j-e_k)$ où $(e_i)$ est la base can. de $V_I$}
\[
F_t\subset V_I\,, \tag*{(60)}
\]
enfin un $k$-module inv. $\underline t\deleted{_E}$ \inserted{$(=\underline t_E)$} et un isom
\[
\unsure{\nu|v}(=\unsure{\nu_E|v_E}):\underline t\deleted{\unsure{_E}}^{\otimes2}\simeq T_{\Sigma_I,t}\,. \tag*{(61)}
\]
Il faut d'abord expliciter calculatoirement $T_{\Sigma_I,t}$
% archive sheet number: 95
\scanpage{97}
comme
\[
\left\{\begin{aligned}
&T_{\Sigma_I,t}\simeq O_t(2)(\underline\omega)\\
&\text{avec }\ O_t(2)\overset{\mathrm{def}}{\simeq}O_t(1)^{\otimes2},\qquad
O_t(1)\simeq V_I/F_t\simeq F_t^\vee(\underline\omega).
\end{aligned}\right. \tag*{(62)}
\]
de sorte que $\unsure{\nu_E|v_E}$ s'interprète comme un isom
\[
\begin{aligned}
\nu:\underline t\deleted{\unsure{_E}}^{\otimes2}&\simeq
\overset{\deleted{\unsure{\vee}}}{F}_t^{\otimes(-2)}(\underline\omega)
\qquad\text{ou aussi comme}\\[-2pt]
&&\check\nu:\underline t^{\otimes-2}\simeq F_t^{\otimes2}(\underline\omega)
\end{aligned}\tag*{(63)}
\]
(NB $\underline\omega=\underline\omega_{\underline I}$, ens des deux ordres circulaires sur $\underline I$). Les données de base sont donc pour nous
\[
\begin{aligned}
[\,I\;(&\mapsto\underline\omega=\underline\omega_I\deleted{\unsure{)}},\deleted{\unsure{F}}\text{ et }V_I=V_I(k)),\\
&F_t\hookrightarrow V_I,\ \underline t,\ \nu:t^{\otimes2}\simeq F_t^{\otimes(-2)}(\underline\omega)\,].
\end{aligned}
\]
Ceci posé, la donnée d'un épinglage de Legendre s'explicite comme un isom
\[
\kappa:\underline t\xrightarrow{\sim}\underline O_t(1)^{\unsure{\sharp|\xi|\natural}} \tag*{(64)}
\]
i.e. comme une famille d'isomorphismes
\[
\left\{\begin{aligned}
&(\kappa_i)_{i\in\mu_4^*},\qquad
\kappa_i:\underline t\xrightarrow{\sim}\underline O_t(1)\simeq F_t^\vee(\underline\omega)\\
&\kappa_{i'}=i\,\kappa_i\qquad\text{si }i'=1/i\text{ est l'autre élément}\\
&\hphantom{\kappa_{i'}=i\,\kappa_i\qquad}\text{de }\mu_4^\times
\end{aligned}\right. \tag*{(65)}
\]
% slanted
\margin{$\mu_4^\times=\mu_4^*(k)$, éléments d'ordre $4$ de $k^*$}
ou encore, d'une famille
\[
\left\{\begin{aligned}
&(\check\kappa_i)_{i\in\mu_4^\times},\qquad \check\kappa_i:\underline t^\vee\simeq F_t(\underline\omega)\\
&\check\kappa_{i'}=\frac1i\,\check\kappa_i\;,
\end{aligned}\right. \tag*{(66)}
\]
satisfaisant à une certaine condition :
\[
\left\{\begin{aligned}
&\kappa_i^{\otimes2}:t^{\otimes2}\xrightarrow{\sim}\overset{\deleted{\unsure{\vee|\wedge}}}{F}_t^{\otimes(-2)}
\quad\text{est de la forme}\\
&\kappa_i^{\otimes2}(\unsure{\delta|\vartheta}^{\otimes2})=\nu(\delta^{\otimes2})\mathbin{\wedge}\omega_i
\end{aligned}\right. \tag*{(67)}
\]
où $\omega_i\in\underline\omega$ est un élément convenable de $\underline\omega$, bien déterminé bien sûr par $\kappa,i$, en fait par $\kappa_i$.
% archive sheet number: 96
\scanpage{98}
On aura alors
\[
\left\{\begin{aligned}
&\deleted{\unsure{\kappa_{i'}=|\omega_{i'}=}}\quad\omega_{i'}=-\omega_i\\
&\text{i.e. }\ i\longmapsto\omega_i\quad\text{est une bijection }\deleted{\unsure{\text{des}|\text{de}}}\\
&\hspace{4em}\alpha=\alpha_\kappa:\mu_4^*\xrightarrow{\sim}\underline\omega
\end{aligned}\right. \tag*{(68)}
\]
% slanted
\margin{\uline{NB} Quand on se donne d'avance l'iso. $\alpha:\mu_4^\times\simeq\underline\omega$, alors les deux relations (67) sont équivalentes \ldots}
On va expliciter la donnée d'un $\kappa$, \uline{associé} à un isom. \deleted{$\alpha:\mu_4^\times\simeq\underline\omega$,} \inserted{donné au préalable} \deleted{\unsure{$\alpha$|à}} $\alpha=(i\longrightarrow\omega_i):\mu_4^\times\simeq\underline\omega$, via les $\check\kappa_i$. On aura, pour $\delta\in\underline t^*$
\[
\check\kappa_i(\check\delta)=\underset{\in F_t}{\underline{\psi_i(\delta)}}\mathbin{\wedge}\omega_i
=\underline{\psi^{\omega_i}(\delta)}\mathbin{\wedge}\omega_i\,, \tag*{-(69)}
\]
où
\[
\bigl(\psi^\omega(\delta)\bigr)_{\omega\in\underline\omega}\ \in (F_t)^{\underline\omega} \tag*{(70)}
\]
est une famille de \deleted{\unsure{éléments|élément}} deux éléments de $F_t$, indexée par $\underline\omega$. La condition\deleted{s} \inserted{(65) ou} (66) s'explicite\deleted{\unsure{nt}} comme
\[
\underbrace{\psi^{\omega_{i'}}(\delta)\mathbin{\wedge}\omega_{i'}}_{
\substack{\psi^{\omega'}(\delta)\mathbin{\wedge}\omega'\\
=-\psi^{\omega'}(\delta)\mathbin{\wedge}\omega\\
\text{si }\deleted{\unsure{\omega'\text{ pour}}}\ \text{\unsure{pos|pose|on pose}}\\
\omega\overset{\mathrm{def}}{=}\omega_i\\
\omega'=\text{autre él. de }\underline\omega}}
=\frac1i\underbrace{\psi^{\omega_i}(\delta)\mathbin{\wedge}\omega_i}_{\psi^\omega(\delta)\mathbin{\wedge}\omega}
\]
i.e. $\psi^{\omega'}(\delta)=\underbrace{-\dfrac1i}_{i}\psi^\omega(\delta)$

i.e
\[
\psi^{\omega'}(\delta)=i_\omega\psi^\omega(\delta)
\qquad\text{où }\ i_\omega\overset{\mathrm{def}}{=}\alpha^{-1}(\omega) \tag*{(71)}
\]
On doit avoir de plus
\[
\psi^\omega(\lambda\delta)=\frac1\lambda\psi^\omega(\delta)\,, \tag*{(72)}
\]
\deleted{qui ex} et enfin il faut exprimer (67),
% archive sheet number: 97
\scanpage{99}
que j'écris sous forme contragrédiente
\[
\check\kappa_i^{\otimes2}(\check\delta^{\otimes2})=
\check\nu(\check\delta^{\otimes2})\mathbin{\wedge}\omega_i\,,
\]
en termes des $\psi^\omega(\delta)$, comme
\[
\psi^\omega(\delta)^{\otimes2}=\check\nu(\check\delta^{\otimes2})\mathbin{\wedge}\omega\,. \tag*{(73)}
\]
\margin{\uline{NB} Moyennant (71) \inserted{(\unsure{resp.|ou} 72)}, les deux relations (73), pour les deux él. $\omega\in\underline\omega$, sont équivalentes, \unsure{resp.|il suffit} \deleted{\unsure{pour}} \inserted{sa} validité pour un $\delta$ \unsure{l'implique|l'impliquent} pour tous.}
On trouve ainsi le

\uline{Lemme}\quad Les épinglages de Legendre pour la donnée $(I,F_t\subset V_I(k),\underline t,\check\nu:\underline{\check t}^{\otimes2}\simeq F_t^{\otimes2}(\underline\omega))$, s'identifient aux systèmes $(\alpha,(\psi^\omega(\delta))_{\omega\in\underline\omega})$, où $\alpha:\mu_4^*\simeq\underline\omega$, et où $(\psi^\omega(\delta))_{\substack{\omega\in\underline\omega\\\delta\in\underline t^*}}\in(F_t)^{\underline\omega\times\underline t^*}$ \deleted{\unsure{satisf|satisfont|satisfaisant}} est une famille indexée par l'ens. \deleted{\unsure{à deux élé|des él.}} \inserted{produit $\underline\omega\times\underline t^*$} \unsure{ments|munts|notons} $\underline\omega=\underline\omega_I$, satisfaisant les conditions (71) (72) (73).

\deleted{\unsure{En}}Nous allons maintenant expliciter la donnée d'une famille à deux indices $(\psi_{\unsure{i|\delta}}^\omega(\delta))_{\substack{\omega\in\underline\omega\\\delta\in\underline t^\times}}$, à val. dans $F_t$, en termes du plongement
\[
F_t\subset V_I=\operatorname{Ker}(k^I\xrightarrow[\text{somme}]{}k)\,,
\]
on voit donc qu'elle s'exprime comme une famille à trois indices d'éléments de $k$ :
\[
\bigl(\psi_i^\omega(\delta)\bigr)_{\substack{i\in I\\\omega\in\underline\omega\\\delta\in\underline t^*}}\,,\qquad
\psi_i^\omega(\delta)\in k \tag*{(74)}
\]
Le fait que la famille d'éléments de $k^I$
% archive sheet number: 98
\scanpage{100}
indexée par $\underline\omega\times\underline t^*$ soit une famille d'éléments de $V_I(k)$ s'exprime par la formule
\[
\sum_{i\in I}\psi_i^\omega(\delta)=0 \tag*{(75)}
\]
les \deleted{\unsure{f}} conditions (71) et (71) s'expriment aussitôt en termes des $\psi_i^\omega(\delta)$, comme
\[
\psi_i^{\omega'}(\delta)=i_\omega\psi_i^\omega(\delta) \tag*{(76)}
\]
\[
\psi_i^\omega(\lambda\delta)=\lambda^{-1}\psi_i^\omega(\delta)\,. \tag*{(77)}
\]
Il reste à exprimer le fait que les éléments $\psi^\omega(\delta)\in V_I(k)$ définis par les $(\omega,\delta)\in\underline\omega\times\underline t\deleted{\unsure{_E}}^*$ sont dans $F_t$ -- moyennant (76) (77) il suffit de l'exprimer pour un seul couple $(\omega,\delta)$ --, et la condition (73). \deleted{On exprime les deux simultanément, en écrivant}
\inserted{Si on définit $F_t$ comme noyau d'une forme linéaire sur $V_I$}
\[
\deleted{\sum_{i\in I}\psi_i^\omega(\delta)e_i-}\tag*{\deleted{(77)}}
\]
\[
F_t=\deleted{\unsure{\sum}}\left\{\sum\lambda_ie_i\mathrel{}\middle|\mathrel{}\sum\lambda_i=0,\ 
\sum c_i\lambda_i\deleted{\unsure{\psi_i^\omega(\delta)|e_i}}=0\right\}
\]
\margin{$(e_i)_{i\in I}$ base canonique de $k^I$}
[\,\uline{NB} l'élément $(c_i)_{i\in I}$ de $\deleted{\unsure{\Sigma}}k^I$ est déterminé modulo multiplication par $k^*$ et addition d'un élément diagonal\,], \deleted{cette}la condition $\psi_{\unsure{\cdot|i}}^\omega(\delta)\in F_t$ s'exprime par la condition calculatoire
\[
\sum_{i\in I}c_i\psi_i^\omega(\delta)=0\qquad . \tag*{(78)}
\]
Il reste \inserted{ensuite} à exprimer (73)
% archive sheet number: 99
\scanpage{101}
\uline{$\overline{\mathrm{IV}}$}\quad Descriptions explicites

par \unsure{fonctions|fonction} \inserted{$\wp$} de Weierstrass
% archive sheet number: 100
\scanpage{102}
\mspage{1}
Soit $E$ courbe elliptique, sur corps alg clos \inserted{$k$} pour commencer (de car $\ne2$ comme d'habitude).

On pose, pour $n\in\mathbb N$
\[
\unsure{\mathcal F_n|\mathscr F_n}=\Gamma(E,n\{0_E\}) \tag*{(1)}
\]
par la formule de RR on a
\[
\operatorname{rg}_k\mathcal F_n=n\qquad\text{si }\ n\geq1 \tag*{(2)}
\]
et
\[
\mathcal F_0=\mathcal F_1=k.1_E\,. \tag*{\unsure{(3)}}
\]
On a \deleted{\unsure{une loi|évidemment}} sur
\[
\mathcal F_\infty\overset{\mathrm{def}}{=}\varinjlim_n\mathcal F_n \tag*{(4)}
\]
une loi de multiplication ass. comm
\[
\mathcal F_n.\mathcal F_m\subset\mathcal F_{n+m}\,, \tag*{(5)}
\]
de sorte que $\mathcal F_*=\coprod\mathcal F_n$ devient une algèbre graduée, dont le $\operatorname{Proj}$ est can. isom. à $E$. On va déterminer $\mathcal F_*$.

Notons la suite exacte
\[
0\mathrel{-}\mathcal F_{n-1}\longrightarrow\mathcal F_n
\xrightarrow{\unsure{\varepsilon_n|\xi_n|\epsilon_n}}t_E^{\otimes n}\longrightarrow0\qquad(n\geq2) \tag*{(6)}
\]
où
\[
t_E\overset{\text{déf}}{=}\underline T_{E,0_E}\,. \tag*{(7)}
\]
On a, pour $\delta\in t_E$, une opération
% archive sheet number: 101
\scanpage{103}
\[
\Theta_\delta:\unsure{\mathcal F_n|\mathscr F_n|F_n}\longrightarrow\mathcal F_{n+1}
\qquad\begin{aligned}&\text{dérivation de degré}\\&\text{$+1$ de l'algèbre}\\&\text{graduée $\mathcal F_*$}\end{aligned}
\tag*{(8)}
\]
qui donne lieu à un hom de suites exactes
\diagram{\begin{tikzcd}[ampersand replacement=\&,column sep=small,row sep=large]
\&0\arrow[r,dash]\&\mathcal F_{n-1}\arrow[r]\arrow[d,"\Theta_\delta"]\&\mathcal F_n\arrow[r]\arrow[d,"\Theta_\delta"]\&t_E^{\otimes n}\arrow[r]\arrow[d,"\mu_\delta"]\&0\\
\text{(9)}\&0\arrow[r]\&\mathcal F_{\deleted{\unsure{n+1|n}}n}\arrow[r]\&\mathcal F_{n+1}\arrow[r]\&t_E^{\otimes n}\arrow[r]\&0
\end{tikzcd}}
où $\mu_\delta$ est la multiplication par $\delta$. De façon plus intrinsèque, \unsure{l'on|on|d'où} a un hom
\[
\begin{aligned}
t_E\otimes\mathcal F_n&\xrightarrow{\delta_n}\mathcal F_{n+1}\\
\delta\otimes\varphi&\longmapsto\Theta_\delta\varphi
\end{aligned}\tag*{(10)}
\]
qui donne lieu à un hom de suites exactes
\diagram{\begin{tikzcd}[ampersand replacement=\&,column sep=small,row sep=large]
\text{(11)}\&0\arrow[r,dash]\&t_E\otimes\mathcal F_{n-1}\arrow[r]\arrow[d,"\delta_{n-1}"]\&\unsure{\mathcal F_E|t_E}\otimes\mathcal F_n\arrow[r]\arrow[d,"\delta_n"]\&t_E\otimes t_E^{\otimes n}\arrow[r]\arrow[d,"{\unsure{\wr|2}\ \mathrm{can}}"]\&0\\
\text{pour }n\geq2\&0\arrow[r,dash]\&\mathcal F_n\arrow[r]\&\mathcal F_{n+1}\arrow[r]\&t_E^{\otimes n+1}\arrow[r]\&0
\end{tikzcd}}
On en conclut que pour $n\geq2$, on a $\operatorname{Ker}\delta_n=\operatorname{Ker}\delta_{n-1}$, donc
\[
\left\{\begin{aligned}
\operatorname{Ker}\delta_n&=\operatorname{Ker}\delta_1=\deleted{k.1_E}\inserted{t_E\otimes}\mathcal F_1=\deleted{k.1_E}t_E\otimes_k1_E\\
\operatorname{Ker}\Theta_\delta&=\mathcal F_1=k.1_E
\end{aligned}\right.\tag*{(12)}
\]
Soit, pour $f\in\mathcal F_n$,
\[
\check f(x)=\check f(-x)\tag*{(13)}
\]
donc
\[
f\longmapsto\check f\qquad \mathcal F_*\longrightarrow\mathcal F_*
\qquad\begin{aligned}&\text{involution de}\\&\text{l'algèbre comm.}\end{aligned}\tag*{(14)}
\]
% archive sheet number: 102
